% GENERATED FILE. Edit canonical lessons, then run npm run generate:books.
% canonical-source-hash: fnv1a64:f78eba63076ec7f1
% canonical-lessons: LA-C187-pax, LA-C187-dux, LA-C187-victoria, LA-C187-libertas, LA-C187-gloria

\chapter{Peace, Leader, Victory, Freedom, Glory}
\label{ch:la-peace-leader-victory-freedom-glory}
\hlchaptermodality{187}

\begin{chapteropening}
\textbf{By the end of this chapter:} \emph{I can say peace, a leader, victory, freedom and glory.}

Complete the last lesson of chapter 187: I can say peace, a leader, victory, freedom and glory.
\end{chapteropening}

\section[pāx, pācis]{pāx, pācis — peace}

\label{lesson:LA-C187-pax}



\begin{quote}
Before the new one: say the Latin for a measure, then the Latin for a pound; a pair of scales.
\end{quote}

\subsection*{You'll want to know: pāx, pācis}
\textbf{pāx, pācis} — \textquotedblleft{}peace\textquotedblright{}.

\begin{cousinweb}
\textbf{pāx} is inherited, usually traced to Proto-Indo-European \emph{*peh2k-}, \textquotedblleft{}to fasten, to fix\textquotedblright{} (EtymDB 2.1): peace as a pact fixed between two sides. \emph{pacīscō}, \textquotedblleft{}to make a bargain\textquotedblright{}, is from the same root (Lewis \& Short).
\end{cousinweb}

\begin{grammarlens}[title={What you've built}]
One more word for this chapter.
\end{grammarlens}

\subsection*{Your turn}
\begin{itemize}
\raggedright
  \item \emph{Say it:} \emph{pāx, pācis}
  \item \emph{Say it:} \emph{pāx, pācis}, once more
  \item \emph{Say it:} say \emph{pāx, pācis}
  \item \emph{Recall:} say the Latin for a measure, then the Latin for a pound; a pair of scales, then say \emph{pāx, pācis} again
\end{itemize}

\begin{tcolorbox}[breakable,colback=teal!4,colframe=teal!35!black,title={Before you move on}]
What does pāx, pācis mean? (\textquotedblleft{}Peace\textquotedblright{}.) Say it once more.
\end{tcolorbox}
\section[dux, ducis]{dux, ducis — a leader}

\label{lesson:LA-C187-dux}



\begin{quote}
Before the new one: say the Latin for a pound; a pair of scales, then the Latin for peace.
\end{quote}

\subsection*{You'll want to know: dux, ducis}
\textbf{dux, ducis} — \textquotedblleft{}a leader\textquotedblright{}.

\begin{cousinweb}
\textbf{dux} is built on \emph{dūcō}, \textquotedblleft{}to lead\textquotedblright{} (Lewis \& Short): the one who leads. English \emph{duke} comes from it through French.
\end{cousinweb}

\begin{grammarlens}[title={What you've built}]
One more word for this chapter.
\end{grammarlens}

\subsection*{Your turn}
\begin{itemize}
\raggedright
  \item \emph{Say it:} \emph{dux, ducis}
  \item \emph{Say it:} \emph{dux, ducis}, once more
  \item \emph{Say it:} say \emph{dux, ducis}
  \item \emph{Recall:} say the Latin for a pound; a pair of scales, then the Latin for peace, then say \emph{dux, ducis} again
\end{itemize}

\begin{tcolorbox}[breakable,colback=teal!4,colframe=teal!35!black,title={Before you move on}]
What does dux, ducis mean? (\textquotedblleft{}A leader\textquotedblright{}.) Say it once more.
\end{tcolorbox}
\section[victōria, victōriae]{victōria, victōriae — victory}

\label{lesson:LA-C187-victoria}



\begin{quote}
Before the new one: say the Latin for peace, then the Latin for a leader.
\end{quote}

\subsection*{You'll want to know: victōria, victōriae}
\textbf{victōria, victōriae} — \textquotedblleft{}victory\textquotedblright{}.

\begin{cousinweb}
\textbf{victōria} is built on \emph{victor}, \textquotedblleft{}a winner\textquotedblright{} (Lewis \& Short), from \emph{vincō}, \textquotedblleft{}to win\textquotedblright{}.
\end{cousinweb}

\begin{grammarlens}[title={What you've built}]
One more word for this chapter.
\end{grammarlens}

\subsection*{Your turn}
\begin{itemize}
\raggedright
  \item \emph{Say it:} \emph{victōria, victōriae}
  \item \emph{Say it:} \emph{victōria, victōriae}, once more
  \item \emph{Say it:} say \emph{victōria, victōriae}
  \item \emph{Recall:} say the Latin for peace, then the Latin for a leader, then say \emph{victōria, victōriae} again
\end{itemize}

\begin{tcolorbox}[breakable,colback=teal!4,colframe=teal!35!black,title={Before you move on}]
What does victōria, victōriae mean? (\textquotedblleft{}Victory\textquotedblright{}.) Say it once more.
\end{tcolorbox}
\section[lībertās, lībertātis]{lībertās, lībertātis — freedom}

\label{lesson:LA-C187-libertas}



\begin{quote}
Before the new one: say the Latin for a leader, then the Latin for victory.
\end{quote}

\subsection*{You'll want to know: lībertās, lībertātis}
\textbf{lībertās, lībertātis} — \textquotedblleft{}freedom\textquotedblright{}.

\begin{cousinweb}
\textbf{lībertās} is built on \emph{līber}, \textquotedblleft{}free\textquotedblright{} (Lewis \& Short), with \emph{-tās}, an ending for a quality. \emph{līber} is usually traced to Proto-Indo-European \emph{*h1lewdheros}, \textquotedblleft{}free\textquotedblright{} (EtymDB 2.1).
\end{cousinweb}

\begin{grammarlens}[title={What you've built}]
One more word for this chapter.
\end{grammarlens}

\subsection*{Your turn}
\begin{itemize}
\raggedright
  \item \emph{Say it:} \emph{lībertās, lībertātis}
  \item \emph{Say it:} \emph{lībertās, lībertātis}, once more
  \item \emph{Say it:} say \emph{lībertās, lībertātis}
  \item \emph{Recall:} say the Latin for a leader, then the Latin for victory, then say \emph{lībertās, lībertātis} again
\end{itemize}

\begin{tcolorbox}[breakable,colback=teal!4,colframe=teal!35!black,title={Before you move on}]
What does lībertās, lībertātis mean? (\textquotedblleft{}Freedom\textquotedblright{}.) Say it once more.
\end{tcolorbox}
\section[glōria, glōriae]{glōria, glōriae — glory}

\label{lesson:LA-C187-gloria}



\begin{quote}
Before the new one: say the Latin for victory, then the Latin for freedom.
\end{quote}

\subsection*{You'll want to know: glōria, glōriae}
\textbf{glōria, glōriae} — \textquotedblleft{}glory\textquotedblright{}.

\begin{cousinweb}
\textbf{glōria} has no sure origin. Lewis \& Short connect it with Greek \emph{kleos}, \textquotedblleft{}fame\textquotedblright{}, and EtymDB 2.1 records a reconstructed older \emph{*gnōria}. Some tie that form to the root of \emph{nōscō}, \textquotedblleft{}to know\textquotedblright{}; neither account is settled.
\end{cousinweb}

\begin{grammarlens}[title={What you've built}]
One more word for this chapter.
\end{grammarlens}

\subsection*{Your turn}
\begin{itemize}
\raggedright
  \item \emph{Say it:} \emph{glōria, glōriae}
  \item \emph{Say it:} \emph{glōria, glōriae}, once more
  \item \emph{Say it:} say \emph{glōria, glōriae}
  \item \emph{Recall:} say the Latin for victory, then the Latin for freedom, then say \emph{glōria, glōriae} again
\end{itemize}

\begin{tcolorbox}[breakable,colback=teal!4,colframe=teal!35!black,title={Before you move on}]
What does glōria, glōriae mean? (\textquotedblleft{}Glory\textquotedblright{}.) Say it once more.
\end{tcolorbox}
